\section{Evidence, Reasoning, and Procedural Limits}

\subsection{The contribution of discovery}

Discovery converted broad disagreement into express admissions and qualified positions.  The defense conceded the physical chronology, qualifying summit location, and later interim control.  The plaintiff conceded missing internal orders, precise operational dates, coordinates, a strike-to-position link, original recordings, and a reporting consensus.  Neither set of concessions decided the proposition.  Together, they defined the evidence on which a jury could decide the remaining inference.

The two technical reports preceded interrogatories and admissions.  Their early formulations therefore provide a point of comparison with the final arguments.  TR-P1's sensitivity analysis initially said that the result changed only if control required permanent annexation or the summit fell within the Golan exclusion.  TR-D1 supplied a third position: temporary control could qualify while the evidence remained insufficient to establish that objective at commencement.  Later plaintiff filings addressed that temporal position through initial position-taking and continuity.  This sequence shows an argument becoming more specific in response to opposition.

The defense also refined its presentation.  TR-D1 contrasted an offensive territorial-control operation with a protective deployment and emphasized the absence of documented combat or an assault plan.  Interrogatory answers and admissions then made explicit that an unopposed operation could qualify, control could be temporary, and defensive motive did not dispose of the claim.  The final defense position used missing operational detail as a limit on evidentiary weight and focused on the timing of the objective.  The development avoided relying on a categorical requirement that the market definition did not contain.

The paper can identify these changes because reports, discovery responses, work notes, and trial filings remain available separately.  A final verdict alone would omit the concessions that made the dispute narrower.  It would also conceal the risk of overstating a conditional admission: the defense's acknowledgment that a class of conduct could qualify was repeatedly used by the plaintiff as support for the conclusion that this instance did qualify.

\subsection{Source diversity and dependence}

The substantive source record comprised two opposing governments' letters, a UN observer statement, and one AP report.  These sources performed different evidentiary functions.  The letters supplied official positions and admissions.  UNDOF supplied reported observations of movement and continued presence, along with security context.  AP supplied details about the summit's location, later fortification, and the retention statement.  Their disagreement about characterization was part of the evidence presented.

The source packages and technical reports increased inspectability without increasing the number of independent accounts of the events.  The defense's independent download of an official letter established access to the same document.  It did not create another account of what happened.  Similarly, the later full PBS capture supported checking the filed excerpt against a complete page retrieved during discovery.  It did not authenticate an original recording of the political statements or determine when the operational decision occurred.

The source rule requires care because the strongest geographical and later-purpose details came from AP.  The plaintiff argued that official confirmation of core acts satisfied the permitted-source route and that AP added corroboration.  The defense accepted the source's publication and credited much of its content while contesting the inference.  The judge's charge permitted source assessment but did not state a detailed rule for combining an official account of deployment with news-only details of later control.  The record supports a description of how the jury used that combination.  It leaves room to question whether the instruction fully resolved the market's source condition.

The initial attachment included Polymarket's Yes resolution and rationale.  Both advocates said that they would use it for the criteria rather than as independent historical proof, and neither offered it as a factual exhibit.  Nevertheless, they had read a supplied answer favoring the plaintiff before researching.  The record supports their stated evidentiary distinction and subsequent use of external sources.  It cannot establish that the supplied outcome had no influence on their framing.

\subsection{Commencement, continuity, and the unit of analysis}

The dominant temporal question was whether the objective inferred from the December 17 posture existed when the operation began around December 7--8.  The plaintiff treated the sequence as one continuous operation.  The defense accepted continuous physical presence but disputed continuous purpose.  This difference persisted from the technical reports through closings and into the votes.

The market period extended through December 31.  A control-seeking operation commenced later within December could therefore raise a separate question even if the initial deployment had a different objective.  Counsel concentrated on the original deployment and whether the later conduct disclosed its original purpose.  Their filed arguments did not develop a distinct theory of a new qualifying offensive commencing on a later date.  The record also lacked the order or precise decision date that would have supported such a separate theory.  This is a limit of the case as litigated, not a finding that a later theory would succeed.

The jury charge used “offensive in method” and “actual authority or dominion” to express the operative criteria.  Those terms gave both sides a way to separate motive, conduct, and control.  The record contains no measured operational threshold for when purposeful military presence becomes an offensive method or when control incidental to a protective deployment becomes an objective to establish territorial authority.  The majority inferred the threshold from the combined conduct.  The dissent found the purpose evidence insufficient at commencement.

The most developed majority explanations gave affirmative reasons for their conclusions.  They named intentional cross-line position-taking, continued occupation, fortification, and conditional withdrawal.  Four jurors, J2, J6, J12, and J15, also relied on the absence of evidence of changed purpose.  That absence can be considered alongside a supported continuity inference, but it cannot alone transfer the burden to the defendant.  J1's dissent and the defense's closing identify the precise risk.  The explanations permit examination of it, while leaving the report without a controlled test of whether a different instruction would change a vote.

\subsection{Vote explanation and access to earlier ballots}

The explanations ranged from a source-by-source and temporal analysis to two-sentence conclusions.  Four plaintiff voters reported high confidence and three reported medium confidence.  The dissent reported medium confidence.  Those labels describe the jurors' own responses.  Neither the model-pool scores nor this case calibrate them as probabilities of historical truth.

The votes were collected sequentially.  The dedicated deliberation prompt supplied the trial transcript and instructions and described round one as having no prior ballot round.  Its prior-ballot section adds earlier-round results only after the first round~\cite{jurorprompt}.  The broader role-visible case response, however, retained public docket entries recording already-filed vote choices and confidence.

This distinction is observable in J2's retained MCP output.  Its full case response includes a docket entry reading “round=1 J1: vote=defendant damages=0.000000 confidence=medium.”  The separate juror-vote and explanation arrays were redacted from the ordinary view.  The full tool response still made the earlier choice available through the docket.  Pi truncated long inline tool results and saved complete responses to files, so availability does not establish that the model read or used the specific entry~\cite[J2 deliberation output and role-visible state]{record}.

The run supplies individual votes under sequential execution with partial access to earlier ballot information.  It supplies no basis for claiming fully isolated ballots or statistically independent judgments.  Shared evidence, instructions, persona, intermediary, and overlapping model families supply additional reasons to distinguish model variety from independence.

\subsection{Judicial activity and the meaning of formal checks}

The judge performed identifiable work: selecting trial mode, entering the pretrial order, screening all thirteen follow-up questions, ruling on the cause challenge, empaneling the jury, settling and delivering instructions, and entering judgment.  The pretrial order preserved the parties' agreements and remaining disputes.  The cause ruling relied on specific words in J10's answer.  The final charge reconciled compatible proposals while leaving the factual inference to the jury.

The proceeding used a documentary trial.  Counsel named public declarants and institutions as sources, but the record contains no live testimony, witness examination, or subpoena-produced internal records.  The two technical reports were prepared by the same model lawyers who advocated the competing outcomes.  Their systematic structure made assumptions and source limits inspectable, but supplied no separate expert witness or independently tested expert method.

Lean's role concerns the encoded procedure.  It enforced opportunity authority, limits, the zero-damages requirement, and the recorded transition sequence.  Its acceptance of an exhibit offer followed an API that lets counsel specify admission status.  Its acceptance of a vote required valid fields and a permitted juror, while the explanation's factual and inferential quality remained a subject for review.

\section{Execution, Cost Records, and Verification}

\subsection{Repeated production opportunities and context growth}

Production generated repeated opportunities to pass on unchanged files.  ADC's \code{updateCase} clears saved opportunity passes, so an accepted production makes remaining optional items available again~\cite[\code{updateCase} and \code{recordOpportunityPass}]{adccore}.  During production, counsel repeatedly encountered files they had already considered.  The plaintiff first passed on the market attachment, then reconsidered the broad request for materials used in its report and produced the attachment.

At 15:00:53, the plaintiff's submitted note explained the changed responsiveness decision:
\begin{quote}
Responsiveness reassessed under RFP 8's broad 'materials considered or relied on' wording: although file-0001 is already shared and was used only for the governing test, it was considered in preparing TR-P1.
\end{quote}
That decision has an evidentiary explanation, but the repeated return of already-passed opportunities has a runtime explanation.  The repeated passes increased interaction without adding new historical sources.

The lawyer sessions completed four context compactions.  Pi reported approximately 923,000--928,000 tokens before those compactions and approximately 15,500--15,700 estimated tokens afterward.  Compaction replaced earlier session content with model-produced summaries and allowed the lawyers to continue.  A further compaction-start event appears at the end of the combined logs without a completed result.

Work-note acknowledgments contributed to the volume of repeated state.  The role API's \code{send\_work\_notes} handling calls \code{opportunityPayloadLocked}, which attaches the role-visible view even when it omits a freshly rendered prompt~\cite{roleapi}.  Thus a short note could receive a growing case record.  Counsel also continued making wait-related model requests while jurors occupied the court turn.  The record does not assign an exact fraction of total token use to either mechanism.

Submitted notes capture source findings, planned concessions, proposed questions, and decisions to pass on unsupported motions.  Other notes announce the next permitted act or repeat a completed assessment.  The notes support reconstruction of strategy, with uneven informational value across entries.

\subsection{Errors by point of failure}

The location of a failure determined which limit and recovery procedure applied.  Pi rejected malformed argument strings before dispatch, so those calls bypassed ADC's decision-attempt counter.  ADC received the overlength theories and nonzero-damages vote, returned a procedural error, and accepted corrected submissions.  A process that ended without its required act used the participant-failure procedure.  That procedure replaced failed candidates during selection and removed a failed sworn juror during deliberation.

The long J14 episode motivated a separate, configurable Pi court-submission error counter.  The later implementation defaults to ten errors per participant opportunity and counts malformed or otherwise rejected court submissions inside Pi.  Court reads and unrelated research tools use their existing handling.  Live tests of that later implementation belong to a separate development record.  They did not alter this case, which retained its original executable and 25-minute deadline.  The report's observations about this run therefore include the full malformed-call interval.

\subsection{Elapsed time and recorded cost}

The unified run lasted approximately two hours and four minutes.  Pleadings, research, discovery, and the pretrial order occupied the first 45 minutes.  Candidate preparation, questionnaires, replacement, follow-up questions, and challenges occupied approximately the next 56 minutes.  The sworn jury was in place at 16:14:04, and judgment followed about 22 minutes later.  The 25-minute malformed-call episode was a substantial part of selection time.

The lawyer logs' token-price estimates sum to \$381.25, including completed compactions.  The lawyers used Codex subscription authentication.  That estimate is a calculation from recorded usage and prices, not an invoice or evidence of an additional charge.  It excludes the court, juror, and separate search requests.

The core's separate provider accounting lists 39 requests, usage observations for 37, and cost observations for 18.  Its observed dollar field totals about \$0.0076.  That incomplete field cannot establish the cost of all court requests.  Juror request records also leave gaps in dollar accounting.

The retained run occupied 444 MiB after completion.  The launcher removed the case's containers, and checks found no retained authentication file or bearer-token file under the run directory.  Persistent case files, workspaces, session records, and process logs remained available for this report.  Disk use by unrelated development caches and other cases falls outside that run-directory measurement.

\subsection{Replay result and reproducibility}

Local certificate verification returned success for 227 transitions.  The \code{adc verify-certificate} command checked the run's \code{case/core} directory through the installed Lean runner with one job, a 6 GiB memory ceiling, 100\% CPU quota, and a 900-second timeout.  Verification replayed the recorded accepted transitions and matched the final state.  The saved state records judgment entered, the proposition demonstrated, seven plaintiff votes, the six-vote requirement, and zero damages.

The database contains 363 event rows, while the certificate contains 227 transitions.  Reads, failed attempts, process-completion records, and other event types explain why those counts measure different things.

The report directory is \path{reports/marxiv/israel-syria-adc/}.  It contains the versioned source and PDF.  The generated run remains outside version control under \path{adc/out/ex04-20260919-04/}.  The full final charge and accepted vote explanations are reproduced here, along with selected discovery and advocacy excerpts.  These excerpts support reading the report independently, while a complete audit of every tool response requires access to the retained run.

\section{Conclusion}

This ADC case produced a declaratory judgment through public-source research, reciprocal discovery, jury selection, documentary trial, instructions, and eight accepted votes.  The lawyers agreed on substantial physical conduct and exposed the missing evidence.  Their dispute focused on offensive method and on the inference from later control to intent at commencement.  Seven jurors accepted the plaintiff's inference, and one gave a reasoned dissent under the same burden.

The record contains examples of procedural work affecting what the factfinder received: the defense obtained the full source page behind an excerpt, both parties made qualified admissions, counsel declined unsupported motions, and a follow-up answer led to a successful cause challenge.  It also records limits: lawyer-controlled admission flags, one principal news report for key later facts, unresolved source-rule detail, access to earlier ballot choices, uneven explanations, and avoidable interaction during waiting and production.

The verified transition sequence establishes completion under the encoded procedure.  The evidence and explanations permit a separate assessment of the resulting judgment.  This single case supports a documented account of that process and disagreement.  Comparisons of accuracy, cost-effectiveness, ballot isolation, or the value of added procedural stages require further cases or controlled comparisons.
